\section{Gabarito e resoluções comentadas — História e Geografia do Maranhão}

\begin{solucao}{08.01}\textbf{Certo.} A França Equinocial foi a tentativa francesa de colonização no Maranhão em 1612; França Antártica refere-se à experiência francesa no Rio de Janeiro no século XVI.\end{solucao}
\begin{solucao}{08.02}\textbf{Certo.} Guaxenduba, em 1614, integra a reação luso-portuguesa contra os franceses.\end{solucao}
\begin{solucao}{08.03}\textbf{Gabarito: B.} A ocupação holandesa de São Luís ocorreu em 1641, depois do ciclo francês do início do século XVII.\end{solucao}
\begin{solucao}{08.04}\textbf{Certo.} A Revolta de Beckman esteve ligada à crise do abastecimento, ao monopólio comercial e a conflitos envolvendo mão de obra indígena e jesuítas.\end{solucao}
\begin{solucao}{08.05}\textbf{Errado.} A adesão maranhense não foi imediata em 1822; consolidou-se em 1823 após conflitos regionais.\end{solucao}
\begin{solucao}{08.06}\textbf{Gabarito: B.} A Balaiada foi revolta regencial de forte base popular, com causas sociais e políticas.\end{solucao}
\begin{solucao}{08.07}\textbf{Certo.} Esses são os limites interestaduais e marítimo do Maranhão.\end{solucao}
\begin{solucao}{08.08}\textbf{Certo.} A distinção aparece expressamente no conteúdo editalício.\end{solucao}
\begin{solucao}{08.09}\textbf{Gabarito: B.} APA é unidade de uso sustentável; parque nacional é unidade de proteção integral.\end{solucao}
\begin{solucao}{08.10}\textbf{Certo.} Babaçu é elemento marcante das formações de cocais.\end{solucao}

\begin{solucao}{08.11}\textbf{Gabarito: A.} La Touche liderou a expedição francesa ligada à França Equinocial e à fundação de São Luís.\end{solucao}
\begin{solucao}{08.12}\textbf{Certo.} Jenipapo ocorreu no Piauí, mas integra o processo regional de consolidação da Independência no norte.\end{solucao}
\begin{solucao}{08.13}\textbf{Certo.} Beckman contestava mecanismos do sistema colonial, não constituía projeto nacional de independência.\end{solucao}
\begin{solucao}{08.14}\textbf{Gabarito: B.} O Vitorinismo deriva da liderança e influência política de Vitorino Freire no século XX.\end{solucao}
\begin{solucao}{08.15}\textbf{Certo.} O episódio ocorreu em contexto de intensa disputa política e mobilização social.\end{solucao}
\begin{solucao}{08.16}\textbf{Gabarito: C.} O Maranhão é espaço de transição ambiental e territorial, o que explica sua diversidade de formações.\end{solucao}
\begin{solucao}{08.17}\textbf{Certo.} Regime de chuvas condiciona rios, vegetação e atividades agropecuárias.\end{solucao}
\begin{solucao}{08.18}\textbf{Certo.} Hidrografia é conteúdo funcional: abastecimento, ocupação, produção e ambiente.\end{solucao}
\begin{solucao}{08.19}\textbf{Gabarito: A.} A Baixada é marcada por planícies e áreas sujeitas a inundações sazonais.\end{solucao}
\begin{solucao}{08.20}\textbf{Certo.} O sul e leste do Maranhão integram a expansão agrícola do MATOPIBA.\end{solucao}
\begin{solucao}{08.21}\textbf{Gabarito: B.} O complexo portuário e as ferrovias são elementos centrais dos corredores de exportação.\end{solucao}
\begin{solucao}{08.22}\textbf{Certo.} São Luís concentra serviços superiores, administração, porto e atividades industriais.\end{solucao}
\begin{solucao}{08.23}\textbf{Certo.} Imperatriz exerce forte centralidade regional no sudoeste maranhense.\end{solucao}
\begin{solucao}{08.24}\textbf{Gabarito: A.} Babaçu = extrativismo vegetal; bumba-meu-boi = cultura; Itaqui = logística.\end{solucao}
\begin{solucao}{08.25}\textbf{Certo.} Urbanização pode coexistir com desigualdade territorial de infraestrutura e serviços.\end{solucao}

\begin{solucao}{08.26}\textbf{Certo.} Guaxenduba pertence ao conflito contra franceses, não à expulsão holandesa.\end{solucao}
\begin{solucao}{08.27}\textbf{Gabarito: B.} A ordem correta é França Equinocial, Guaxenduba, ocupação holandesa e Beckman.\end{solucao}
\begin{solucao}{08.28}\textbf{Certo.} A Balaiada combinou disputas políticas e profundas tensões sociais, com ampla participação popular.\end{solucao}
\begin{solucao}{08.29}\textbf{Certo.} A Revolução de 1930 reorganizou a relação entre centro e estados e alterou os grupos políticos locais.\end{solucao}
\begin{solucao}{08.30}\textbf{Gabarito: A.} Os eventos atravessam disputa colonial, crise do sistema colonial e Regência.\end{solucao}
\begin{solucao}{08.31}\textbf{Certo.} A localização de transição ajuda a explicar gradientes de chuva e mosaico vegetal.\end{solucao}
\begin{solucao}{08.32}\textbf{Gabarito: A.} Babaçuais são característicos da Mata dos Cocais.\end{solucao}
\begin{solucao}{08.33}\textbf{Certo.} Manguezais são ecossistemas costeiro-estuarinos.\end{solucao}
\begin{solucao}{08.34}\textbf{Gabarito: B.} Essa é a distinção central entre as duas categorias de unidade de conservação.\end{solucao}
\begin{solucao}{08.35}\textbf{Certo.} Planejamento territorial deve integrar fatores físicos e riscos ambientais.\end{solucao}
\begin{solucao}{08.36}\textbf{Gabarito: A.} A EFC integra o corredor de Carajás até a região portuária de São Luís.\end{solucao}
\begin{solucao}{08.37}\textbf{Certo.} A estrutura econômica é diversificada e territorialmente desigual.\end{solucao}
\begin{solucao}{08.38}\textbf{Certo.} Crescimento econômico não implica automaticamente redução de desigualdades.\end{solucao}
\begin{solucao}{08.39}\textbf{Gabarito: A.} São Luís e Imperatriz são centros regionais com funções diferentes e áreas de influência próprias.\end{solucao}
\begin{solucao}{08.40}\textbf{Certo.} As manifestações culturais resultam de matrizes africanas, indígenas e europeias articuladas historicamente.\end{solucao}

\begin{solucao}{08.41}\textbf{Gabarito: A.} Atividade agrícola depende de condições naturais, tecnologia, terra, crédito e acesso logístico aos mercados.\end{solucao}
\begin{solucao}{08.42}\textbf{Certo.} Uso da terra é tema integrado de geografia física, econômica e ambiental.\end{solucao}
\begin{solucao}{08.43}\textbf{Gabarito: B.} Planejamento deve considerar dinâmica de cheias, ocupação e drenagem, reduzindo exposição e vulnerabilidade.\end{solucao}
\begin{solucao}{08.44}\textbf{Certo.} Infraestrutura logística pode gerar ganhos concentrados e não assegura distribuição territorial uniforme.\end{solucao}
\begin{solucao}{08.45}\textbf{Gabarito: A.} A sequência atravessa Colônia, Império e República.\end{solucao}
\begin{solucao}{08.46}\textbf{Certo.} A consolidação da Independência ocorreu em escala regional e envolveu províncias vizinhas.\end{solucao}
\begin{solucao}{08.47}\textbf{Gabarito: A.} Cultura é expressão espacial e histórica das relações sociais e identidades coletivas.\end{solucao}
\begin{solucao}{08.48}\textbf{Certo.} A distinção permite ler tanto limites estaduais quanto drenagem interna.\end{solucao}
\begin{solucao}{08.49}\textbf{Gabarito: A.} Desenvolvimento sustentável exige visão multidimensional do território.\end{solucao}
\begin{solucao}{08.50}\textbf{Certo.} Essa é a síntese adequada para interpretar o estado em perspectiva histórica e geográfica integrada.\end{solucao}
